\section{Equality of the two infima}\label{sec:infima}

The obstruction above concerns attainment, not the value of the Monge
infimum. We prove this directly for the constructed marginal. The
partition and cost-oscillation argument adapts the approximation method
of Colombo and Di Marino~\cite[author version, Lemma~2.3 and proof of
Theorem~2.1]{CDM2015}. Their general theorem treats the Coulomb singularity
and cyclic maps for identical nonatomic marginals. Here the separated
components allow an elementary proof, including the required Borel maps.

\begin{proposition}\label{prop:equal-infima}
For the probability measure \(\mu\) of Theorem~\ref{thm:main},
\[
 \inf_{\substack{T_2,T_3:\R^3\to\R^3\ \mathrm{Borel}\\
                 (T_2)_\#\mu=(T_3)_\#\mu=\mu}}
 \int c(x,T_2(x),T_3(x))\dd\mu(x)=\E(\mu).
\]
Every pair in this infimum has cost strictly greater than the common
infimum, but pairs with finite cost approach it.
\end{proposition}

The approximation keeps the finite component structure of the optimizer
\(\pi_0\). We partition the components into small cells and prescribe the
mass of a graph plan in each product of three cells to equal that of
\(\pi_0\). Within each product the two plans may differ, but the components
remain separated and the cost oscillation tends to zero with the cell
diameter. To realize the prescribed masses, we split each source cell
into pieces and send each piece into two specified target cells, with
the normalized restriction of \(\mu\) as each target law.
The following elementary transport fact supplies those maps; its proof
is given at the end of the section.

\begin{lemma}\label{lem:borel-transport}
Let \(d\ge2\) be an integer. If \(\alpha\) and \(\beta\) are absolutely
continuous Borel probability measures on \(\R^d\), there is a Borel map
\(Q:\R^d\to\R^d\) such that \(Q_\#\alpha=\beta\).
\end{lemma}

\begin{proof}[Proof of Proposition~\ref{prop:equal-infima}]
Write \(K_0=\supp g\), and let \(K_1,\ldots,K_4\) be its images under
the four maps in \(\calH\). These are nonempty compact sets in the five
disjoint components, and \(\mu(K)=1\) for \(K=\bigcup_{a=0}^4K_a\).
In particular,
\[
 \eta=\min_{a\ne b}\operatorname{dist}(K_a,K_b)>0.
\]
The optimal coupling \(\pi_0\) constructed in Section~\ref{sec:proof}
is supported on products of three distinct component sets.

Fix \(\delta>0\). Partition each \(K_a\) into finitely many Borel cells
of diameter at most \(\delta\), for example by a finite grid of half-open
cubes. Discard the cells of zero \(\mu\)-mass and denote the remaining
cells by \(A_1,\ldots,A_N\). They partition a set of full \(\mu\)-measure
and each belongs to a single \(K_a\). Put
\[
 q_{ijk}=\pi_0(A_i\times A_j\times A_k).
\]
The marginal constraints on \(\pi_0\) read
\[
 \sum_{j,k}q_{ijk}=\mu(A_i),\qquad
 \sum_{i,k}q_{ijk}=\mu(A_j),\qquad
 \sum_{i,j}q_{ijk}=\mu(A_k).
\]
For each fixed \(i\), split \(A_i\) into Borel sets \(B_{ijk}\) of masses
\(q_{ijk}\), ignoring zero-mass pieces. Such a split can be made by
successive first-coordinate slabs: the first-coordinate marginal of
\(\mu|_{A_i}\) is absolutely continuous, so its cumulative distribution
is continuous and takes every value between zero and \(\mu(A_i)\).
Endpoints of the slabs have zero mass; any unused null set is included
in a fixed piece.

For each \(q_{ijk}>0\), Lemma~\ref{lem:borel-transport} supplies maps on
\(B_{ijk}\) sending the probability \(q_{ijk}^{-1}\mu|_{B_{ijk}}\),
respectively, to
\[
 \frac{\mu|_{A_j}}{\mu(A_j)}
 \quad\hbox{and}\quad
 \frac{\mu|_{A_k}}{\mu(A_k)}.
\]
Reset any exceptional values outside the indicated target cell to a
fixed point in that cell. These changes are Borel and occur on null
sets. Glue the maps over the finitely many pieces to obtain \(T_2,T_3\),
and give them a fixed value on the remaining Borel null set in \(\R^3\).
Different source pieces may use the same normalized target law. Their
weights \(q_{ijk}\) add to the full mass of the target cell, so no target
partition is needed. Indeed,
for any Borel set \(D\),
\[
 \mu(T_2^{-1}(D))
 =\sum_{i,j,k:q_{ijk}>0}q_{ijk}
       \frac{\mu(D\cap A_j)}{\mu(A_j)}
 =\sum_j\mu(D\cap A_j)=\mu(D),
\]
and the same calculation proves \((T_3)_\#\mu=\mu\).
The graph coupling \(\pi_\delta=(\Id,T_2,T_3)_\#\mu\) gives each product
\(A_i\times A_j\times A_k\) exactly the mass \(q_{ijk}\).

Whenever \(q_{ijk}>0\), the three cells lie in three distinct component
sets. Thus all pairwise distances in their product are at least
\(\eta\), including after the redistribution within the cells. If
\(x\) and \(y\) are two triples in the same such product, then
\(|x_\ell-y_\ell|\le\delta\) for each coordinate. Consequently
\[
 \left|\frac1{|x_\ell-x_r|}-\frac1{|y_\ell-y_r|}\right|
 \le\frac{2\delta}{\eta^2},
 \qquad |c(x)-c(y)|\le\frac{6\delta}{\eta^2}.
\]
Both costs are finite, bounded by \(3/\eta\) on all active products.
Since \(\pi_\delta\) and \(\pi_0\) have identical masses on these
products, summing their oscillation bounds gives
\[
 \left|\int c\dd\pi_\delta-\int c\dd\pi_0\right|
 \le\frac{6\delta}{\eta^2}.
\]
Letting \(\delta\downarrow0\) proves the required upper bound on the
Monge infimum. The reverse bound follows because every admissible graph
plan is a coupling. Strict nonattainment is Theorem~\ref{thm:main}.
\end{proof}

The argument uses a uniform separation estimate on the approximating
plans themselves; it does not pass the singular Coulomb cost through
weak convergence. No classification or uniqueness of optimal couplings
is needed.

\begin{proof}[Proof of Lemma~\ref{lem:borel-transport}]
We construct a Borel map from each such probability to uniform measure
on \((0,1)^d\), and a Borel map in the reverse direction. Let \(f\) be a
nonnegative Borel probability density. The forward map is Rosenblatt's
successive conditional-distribution transform~\cite{Rosenblatt1952}; we
also give the inverse quantile construction and make the Borel choices
explicit in both directions.

\emph{The forward map.}
We successively transform each coordinate by its conditional distribution
function given the preceding coordinates. Write
\(x_{1:j}=(x_1,\ldots,x_j)\), with \(x_{1:0}\) the empty prefix, and set
\[
 f_{1:j}(x_{1:j})=\int_{\R^{d-j}}f(x_{1:j},z)\dd z
 \quad(1\le j<d),\qquad f_{1:d}=f,\qquad f_{1:0}=1.
\]
The prefix densities are Borel, have integral one, and are finite almost
everywhere. Tonelli's theorem gives
\(\int_\R f_{1:j}(a,r)\dd r=f_{1:j-1}(a)\) for every prefix \(a\),
including the empty prefix when \(j=1\). For \(j=1,\ldots,d\), define
\[
 F_j(a,t)=\frac{\int_{-\infty}^t f_{1:j}(a,r)\dd r}
                 {f_{1:j-1}(a)},\qquad a\in\R^{j-1}.
\]
This formula is used where its denominator belongs to \((0,\infty)\);
on the remaining fibers use instead the fixed distribution function
\(t\mapsto(1+e^{-t})^{-1}\). Thus every fiber is a continuous distribution
function with limits zero and one, and all \(F_j\) are jointly Borel.
For \(j\ge2\), the exceptional prefix fibers have zero probability under
\(f_{1:j-1}(a)\dd a\); there are no exceptional fibers when \(j=1\).

For a random variable with a continuous distribution function \(F\),
the variable \(F(X)\) is uniform on \((0,1)\). Indeed, continuity and the
limits of \(F\) allow a cut at each probability level \(u\in(0,1)\);
any interval on which \(F=u\) has probability zero, and hence
\(\mathbb P(F(X)\le u)=u\). At stage \(j\), this identity and Fubini show
that, conditionally on the preceding original prefix, the new transformed
coordinate is uniform.
It is therefore independent of that prefix, and hence of all preceding
transformed coordinates. A finite induction over \(j=1,\ldots,d\) shows that the Borel map
\[
 R_f(x)=\bigl(F_j(x_{1:j-1},x_j)\bigr)_{j=1}^d
\]
sends \(f\dd x\) to product uniform measure on \((0,1)^d\). Its values lie
in \([0,1]^d\), including on exceptional fibers. This argument does not
require a density to be positive everywhere or a distribution function to
be strictly increasing.

\emph{The reverse map.}
Given independent uniform coordinates, we successively recover the
original conditional laws by quantiles. For each \(j\)
and \(0<u<1\), define
\[
 q_j(a,u)=\inf\{r\in\mathbb Q:F_j(a,r)\ge u\}.
\]
This value is finite, and
\[
 \{(a,u):q_j(a,u)<t\}
 =\bigcup_{\substack{r\in\mathbb Q\\r<t}}
      \{(a,u):F_j(a,r)\ge u\}
\]
proves joint Borel measurability on \(0<u<1\). At \(u=0,1\), define
\(q_j(a,u)=0\). For \(0<u<1\), continuity of \(F_j(a,\cdot)\) gives
\[
 q_j(a,u)\le t\quad\Longleftrightarrow\quad u\le F_j(a,t).
\]
Consequently \(q_j(a,U)\) has distribution function \(F_j(a,\cdot)\)
for uniform \(U\).
For \(u\in[0,1]^d\), recursively set
\[
 z_j=q_j(z_{1:j-1},u_j)\quad(j=1,\ldots,d),
 \qquad S_f(u)=(z_1,\ldots,z_d),
\]
starting with the empty prefix. This finite recursion defines a Borel map.
For independent uniform inputs, induction shows that after step \(j\) the
output prefix has density \(f_{1:j}\): the preceding prefix has density
\(f_{1:j-1}\), its exceptional fibers have probability zero, and on every
other fiber the quantile has the conditional distribution just defined.
Thus \(S_f\) has output law \(f\dd x\).
If \(f\) and \(g\) are densities of \(\alpha\) and \(\beta\), respectively,
then \(S_g\circ R_f\) has the required property.
\end{proof}
